\begin{abstract}
Large-scale unsupervised language models (LMs) acquire extensive general knowledge and certain reasoning abilities; however, their behavior remains challenging to direct with precision due to their entirely unsupervised training regimes. Current approaches to improving model steerability typically involve gathering human-annotated pairwise comparisons of model outputs and then fine-tuning the unsupervised LM to align with these preferences, most commonly through reinforcement learning from human feedback (RLHF). Yet, RLHF entails a multi-stage and often unstable process, requiring the construction of a reward model that captures human preference judgments, followed by reinforcement learning-based adaptation of the LM to maximize the surrogate reward while constraining deviation from the original parameters. This paper presents a novel reward model parameterization within the RLHF framework, which enables one to analytically derive the associated optimal policy, thereby reducing the standard RLHF objective to a straightforward classification loss. Our proposed method, termed \textit{Direct Preference Optimization} (DPO), is both reliable and effective, offering a lightweight alternative that obviates the necessity for generation sampling or extensive hyperparameter search during fine-tuning. Empirical results indicate that DPO fine-tunes LMs to adhere to human preferences at least as well as, or superior to, established techniques. In particular, DPO fine-tuning surpasses PPO-based RLHF in sentiment steering, and achieves equal or improved quality in summarization and single-turn dialogue tasks, all while being significantly easier to apply and optimize.
\end{abstract}

\section{Introduction}
Large-scale unsupervised language models (LMs), when trained on extensive corpora, exhibit a range of unexpected capabilities~\citep{chowdhery2022palm, brown2020language, touvron2023llama,bubeck2023sparks}. Nevertheless, these models are constructed from human-generated data that reflects a diverse array of intentions, skill levels, and priorities. Not all of these are appropriate for replication: for instance, although it is valuable for an AI coding assistant to \textit{recognize} common programming errors to provide corrections, we ultimately desire that the system’s code generation favors the subset of high-caliber programming skills, even if those examples are infrequent in the training set. Likewise, while it may be useful for a language model to \textit{identify} a misconception held by half the population, we emphatically do not want it to assert this misconception as truth in half of relevant responses. Thus, the ability to \emph{specify and elicit desired responses and behaviors}—drawing selectively from the model's extensive \textit{knowledge and skill} repertoire—is essential for the development of AI systems that are reliable, high-performing, and manageable \citep{ouyang2022training}. Common practice involves aligning language models with human values using reinforcement learning (RL), yet, as we will demonstrate, the conventional RL-based approach to this task can actually be solved exactly through a straightforward binary cross-entropy loss, leading to significant streamlining of the preference learning process.

Broadly, established approaches incorporate preferred behaviors into language models by leveraging carefully constructed datasets that reflect human evaluations of what is beneficial or secure. This phase of learning from preferences succeeds a preliminary, large-scale unsupervised pre-training step conducted over vast text resources. While the most direct method of preference learning consists of supervised fine-tuning on exemplars of superior responses provided by humans, the most effective set of techniques so far has been reinforcement learning from human or AI feedback (RLHF/RLAIF; \citep{christiano2017deep,bai2022constitutional}). In RLHF, a reward model is fitted to annotated preference data and, subsequently, reinforcement learning is used to adapt the policy of the language model to prefer outputs rated as highly rewarding, all while limiting divergence from the original model. Despite the strong results achieved with RLHF for dialogue and code generation, the RLHF methodology remains more intricate than standard supervised approaches, requiring the training of several LMs as well as sampling from the evolving policy during training, which results in considerable computational overhead.

This work introduces a novel approach for aligning language models with human preferences without employing explicit reward modeling or reinforcement learning frameworks. Specifically, we introduce \textit{{Direct Preference Optimization} (DPO)}, an algorithm that implicitly targets the same objective function as popular RLHF techniques—namely, reward maximization under a KL-divergence regularization—yet features a streamlined implementation and training regime. Conceptually, DPO adjusts the model by boosting the log likelihood of preferred responses over non-preferred ones; crucially, it utilizes adaptive, instance-level importance weights to counteract issues such as model collapse that can arise when applying straightforward probability ratio-based objectives. DPO, akin to prior approaches, is grounded in a formal preference model (e.g., the Bradley-Terry model; \cite{bradley1952rankanalysis}) to quantify the correspondence between a reward function and observed preference data. In contrast to standard pipelines—which train a separate reward model via a preference loss and subsequently train a policy to optimize that learned reward—DPO reparameterizes the preference loss so that it directly governs the policy parameters. Leveraging collections of human-annotated preferences over model outputs, DPO can be applied to optimize the policy via a binary cross entropy objective, effectively inducing the policy corresponding to an implicit reward that best fits the data.

The principal contribution of this paper is Direct Preference Optimization (DPO): a straightforward method for preference-driven language model training that avoids reinforcement learning. Experimental results demonstrate that DPO matches or surpasses leading alternatives—including RLHF using PPO—in several preference learning benchmarks, such as sentiment conditioning, summarization, and conversational tasks, with models up to 6B parameters.

Self-supervised language models of greater size have demonstrated the capacity to perform certain tasks in a zero-shot setting \citep{radford2019language}, as well as with limited examples provided in prompts \citep{gpt3,megatron,chowdhery2022palm}. Nevertheless, their effectiveness on downstream tasks and alignment with user objectives is substantially enhanced through fine-tuning on collections of instructions and corresponding human-authored completions \citep{mishra-etal-2022-cross,sanh2022multitask,chung2022scaling,thoppilan2022lamda}. This process, termed `instruction-tuning,’ permits large language models (LLMs) to generalize to novel instructions beyond those present in the tuning set, thereby generally improving model usability \citep{chung2022scaling}. In spite of the efficacy of instruction tuning, it is often less resource-intensive to obtain \textit{relative} human assessments of response quality than expert-labeled examples. As a result, subsequent research has explored fine-tuning LLMs using datasets constructed from human preferences, resulting in performance gains in tasks like translation \citep{kreutzer-etal-2018-reliability}, summarization \citep{stiennon2022learning,ziegler2020finetuning}, story-telling \citep{ziegler2020finetuning}, and instruction-following \citep{ouyang2022training,ramamurthy2023is}. The standard methodology involves first training a neural reward function to reflect preferences encoded by a preference model (such as the Bradley-Terry model \citep{bradley1952rankanalysis}), and subsequently fine-tuning the language model with reinforcement learning—using algorithms like REINFORCE \citep{williams1992reinforce}, proximal policy optimization (PPO; \cite{schulman2017proximal}), or related techniques \citep{ramamurthy2023is}—to optimize this learned reward. A related research direction utilizes LLMs, already tuned to follow instructions and trained with human feedback, to autonomously generate new synthetic preference data, targeting qualities such as safety or non-harmfulness \citep{bai2022constitutional}. In these cases, human guidance may be restricted to the provision of a rubric for model-based annotation. Collectively, these approaches synthesize two research trajectories: one exploring the use of reinforcement learning to train language models across various objectives~\citep{Ranzato2015SequenceLT,paulus2018a,wu2018learning}, and the other focused on more general frameworks for preference-based learning from humans \citep{christiano2017deep,kupcsik2018learning}. Notwithstanding the advantages of leveraging relative preferences, the application of reinforcement learning for fine-tuning large language models introduces substantial practical barriers; the present work addresses this by introducing a theoretically-grounded method to optimize relative preference data without the necessity of RL.

Research on learning policies from preferences extends beyond language-based tasks to both bandit and reinforcement learning frameworks, leading to the development of a variety of methods. Within the bandit context, learning from preferences or rankings instead of explicit rewards is referred to as the contextual dueling bandit (CDB; \cite{yue2012karmed,dudik2015contextual}). When absolute reward signals are absent, the theory underlying CDBs replaces the standard optimal policy concept with that of a \textit{von Neumann winner}—a policy achieving an expected win rate of at least 50\% when compared to any other candidate policy \citep{dudik2015contextual}. Nevertheless, CDB algorithms typically involve collecting preference labels interactively in an online setting, while in the domain of learning from human preferences, models are more often trained using a fixed dataset of offline, preference-annotated action pairs \citep{yan2022human}. In parallel, \textit{preference-based RL} (PbRL) refers to methods that optimize policies based on binary preferences produced by an unknown ‘scoring’ function rather than by reward signals \citep{BusaFekete2014,ruiz2023dueling}. Numerous algorithms address PbRL, including those that make use of off-policy preference data, though they usually entail first estimating the underlying scoring (reward) function and then performing optimization on it \citep{jain2013learning,BusaFekete2014,christiano2017deep,sadigh2017active,kupcsik2018learning}. In contrast to this two-step paradigm, our approach introduces a unified stage of policy optimization that seeks to directly maximize preference satisfaction.

\section{Preliminaries}\label{section:prelims}

Here, we provide an overview of the RLHF pipeline as outlined in \citeauthor{ziegler2020finetuning} (see also \citep{stiennon2022learning, bai2022training, ouyang2022training}), which is generally organized into three main components: (1) supervised fine-tuning (SFT); (2) collection of preference data and reward modeling; and (3) reinforcement learning optimization.

\textbf{SFT}: The process of RLHF customarily begins with supervised fine-tuning of a large language model (LM), using labeled data that reflects high-quality completions for the tasks at hand (e.g., dialogue, summarization). This produces an initial model, denoted as $\pi^\text{SFT}$.

\textbf{Reward Modeling Phase}: Following SFT, the model $\pi^\text{SFT}$ is used to generate candidate response pairs $(y_1, y_2) \sim \pi^\text{SFT}(y \mid x)$, conditioned on a prompt $x$. Human annotators are then tasked with expressing a preference between these responses, indicated as $y_w\succ y_l \mid x$, with $y_w$ and $y_l$ representing the selected and rejected completions among $(y_1, y_2)$, respectively. These observed preferences are assumed to result from an unobserved, latent reward function $r^*(y, x)$. Various preference modeling techniques have been explored for this purpose, with the Bradley-Terry (BT) \cite{bradley1952rankanalysis} model frequently employed. The framework also permits alternative ranking models, such as the Plackett-Luce model \citep{plackett1975analysis, luce2012individual}, if more comprehensive ranking data is available. In the BT approach, the probability of the observed human preference $p^*$ is expressed as follows:

\begin{equation}\label{eq:bradley-terry}
    p^*(y_1\succ y_2 \mid x)=\frac{\exp\left(r^*(x, y_1)\right)}{\exp\left(r^*(x, y_1)\right) + \exp\left(r^*(x, y_2)\right)}.
\end{equation}

Suppose we are provided with a fixed dataset of preference comparisons, $\mathcal{D}=\bigl\{x^{(i)}, y_w^{(i)}, y_l^{(i)}\bigr\}_{i=1}^N$, collected according to $p^*$. We can introduce a reward model $r_{\phi}(x, y)$ with learnable parameters and fit it by maximizing the likelihood of the observed preferences. Treating this as a binary classification task, the associated negative log-likelihood loss is given by

\begin{equation}\label{eq:reward_model}
    \mathcal{L}_R(r_{\phi}, \mathcal{D}) = -\mathbb{E}_{(x, y_w, y_l)\sim \mathcal{D}}\bigl[\log \sigma(r_{\phi}(x, y_w)- r_{\phi}(x, y_l))\bigr]
\end{equation}

where $\sigma$ denotes the logistic sigmoid. For language models, $r_{\phi}(x, y)$ is commonly initialized with the SFT model $\pi^\text{SFT}(y \mid x)$, augmented by a linear transformation on top of the transformer's output to generate a scalar reward value~\cite{ziegler2020finetuning}. To ensure stability and control variance, rewards are often normalized so that $\mathbb{E}_{x,y\sim \mathcal{D}}\left[r_\phi(x, y)\right] = 0$ for any given $x$.

\textbf{RL Fine-Tuning Phase}: In the reinforcement learning phase, the acquired reward model serves as a signal to refine the language model further. Consistent with previous research~\citep{jaques2017sequence, jaques2020human}, the objective for optimization takes the form

\begin{equation}\label{eq:RL}
\max_{\pi_{\theta}}  \mathbb{E}_{x\sim \mathcal{D}, y\sim \pi_{\theta}(y \mid x)}\bigl[r_{\phi}(x, y)\bigr] - \beta\mathbb{D}_{\textrm{KL}}\bigl[\pi_{\theta}(y\mid x)\mid \mid \pi_\text{ref}(y\mid x)\bigr],
\end{equation}

where $\beta$ regulates how much $\pi_{\theta}$ is allowed to diverge from a reference policy $\pi_\text{ref}$—typically, the base SFT model $\pi^\text{SFT}$. Generally, the policy $\pi_\theta$ is also initialized from $\pi^\text{SFT}$. This regularization is crucial: it constrains the updated model within the regions where the reward predictions are reliable and discourages mode collapse by sustaining generation diversity. Since generating text is a discrete process, gradients do not flow through this objective directly, so reinforcement learning approaches are employed instead. The prevailing strategy~\cite{ziegler2020finetuning, stiennon2022learning, bai2022training, ouyang2022training} involves defining a combined reward function ${r(x, y) = r_{\phi}(x, y) -\beta (\log \pi_{\theta}(y\mid x) - \log \pi_\text{ref}(y\mid x))}$ and optimizing this target using PPO~\cite{schulman2017proximal}.

\section{Direct Preference Optimization}\label{sec:DPO}

Recognizing the practical challenges encountered when employing RL-based algorithms for large-scale tasks such as language model fine-tuning, we seek an alternative policy optimization framework grounded directly in preference data. In contrast to traditional RLHF pipelines that first train a reward model and subsequently optimize it with RL, we exploit a particular reward parameterization that permits analytic derivation of its optimal policy, obviating the need for an explicit RL loop. As discussed in detail below, our main contribution is to establish a closed-form correspondence between reward functions and their optimal policies, which facilitates converting loss objectives from the reward space into the policy space. This reparameterization circumvents the requirement to train a distinct reward model, but still permits optimization within established frameworks of human preference—such as the Bradley-Terry model. Effectively, this means the policy network not only generates text but simultaneously enc

\textbf{DPO Objective Derivation.} To begin, we adopt the RL objective specified in Eq.~\ref{eq:RL} for a general reward function $r$, consistent with earlier studies. Prior literature~\citep{peters2007reinforcement, peng2019advantage, korbak2022reinforcement, go2023aligning} shows that the solution maximizing the reward with a KL penalty, as given by Eq.~\ref{eq:RL}, leads to an optimal policy of the form:

\begin{equation}\label{eq:op_policy}
    \pi_r(y\mid x) = \frac{1}{Z(x)}\pi_\text{ref}(y\mid x)\exp\left(\frac{1}{\beta}r(x, y)\right),
\end{equation}

Here, the normalization constant $Z(x) = \sum_{y}\pi_\text{ref}(y\mid x)\exp\left(\frac{1}{\beta}r(x, y)\right)$ is known as the partition function. The derivation is provided in Appendix \ref{app:derivation1}. Even when $r$ is estimated via MLE by $r_{\phi}$ as a proxy for the true reward $r^*$, evaluating $Z(x)$ is still computationally intensive~\citep{korbak2022reinforcement, go2023aligning}, creating obstacles for this parameterization in practice. Nevertheless, by algebraically manipulating Eq.~\ref{eq:op_policy}, we can rewrite the reward as a function of its corresponding optimal policy $\pi_r$, the baseline policy $\pi_\text{ref}$, and the partition function $Z(\cdot)$. Specifically, taking logarithms of both sides of Eq.~\ref{eq:op_policy} yields:

\begin{equation}\label{eq:main_eq}
    r(x,y) =\beta \log \frac{\pi_r(y\mid x)}{\pi_\text{ref}(y\mid x)} + \beta \log Z(x).
\end{equation}

This reparameterization can be employed for the ground-truth reward $r^*$ and the associated optimal policy $\pi^*$. Importantly, since the Bradley-Terry model is sensitive only to differences between two rewards—${p^*(y_1 \succ y_2 \mid x) = \sigma(r^*(x, y_1) - r^*(x, y_2))}$—substituting Eq.~\ref{eq:main_eq} for $r^*(x,y)$ in the preference model Eq.~\ref{eq:bradley-terry} causes the partition function to cancel out. Thus, the probability of human preference may be framed exclusively in terms of the optimal policy $\pi^*$ and reference policy $\pi_\text{ref}$. The resulting RLHF-optimal policy $\pi^*$ under the Bradley-Terry framework satisfies:

\begin{equation}\label{eq:objective}
    p^*(y_1\succ y_2 \mid x)=\frac{1}{1 + \exp\left(\beta \log \frac{\pi^*(y_2\mid x)}{\pi_\text{ref}(y_2\mid x)} - \beta \log \frac{\pi^*(y_1\mid x)}{\pi_\text{ref}(y_1\mid x)}\right)}
\end{equation}

For a full derivation, refer to Appendix~\ref{app:derivation2}. Although the discussion above centers on the Bradley-Terry model, analogous formulations for more general Plackett-Luce models~\citep{plackett1975analysis, luce2012individual} are provided in Appendix~\ref{app:plackett_luce_models}.

Since we have reformulated the preference probability directly in terms of the optimal policy, this enables us to define a maximum likelihood training objective for a parametrized policy $\pi_\theta$. Mirroring the approach used for reward modeling (see Eq.~\ref{eq:reward_model}), our policy optimization objective becomes:

\begin{equation}\label{eq:optimum_model}
    \mathcal{L}_\text{DPO}(\pi_{\theta}; \pi_\text{ref}) = -\mathbb{E}_{(x, y_w, y_l)\sim \mathcal{D}}\left[\log \sigma \left(\beta \log \frac{\pi_{\theta}(y_w\mid x)}{\pi_\text{ref}(y_w\mid x)} - \beta \log \frac{\pi_{\theta}(y_l\mid x)}{\pi_\text{ref}(

In this approach, we infer an implicit reward function through an alternative parametrization, for which the optimal policy corresponds directly to $\pi_\theta$. Furthermore, because this method is mathematically analogous to fitting a reparameterized Bradley-Terry model, it benefits from certain theoretical guarantees—most notably, consistency under suitable assumptions regarding the distribution of preference data \cite{bong2022generalized}. Section~\ref{sec:theory} expands on the theoretical analysis of DPO and situates it relative to related research.

\textbf{What is the effect of the DPO update?} To elucidate the mechanism behind DPO, one can investigate the gradient of its loss function, $\mathcal{L}_\text{DPO}$. The gradient with respect to model parameters $\theta$ is given by:

\begin{multline*}\label{eq:gradient}
    \nabla_\theta \mathcal{L}_\text{DPO}(\pi_\theta;\pi_\text{ref}) = \\ -\beta\mathbb{E}_{(x, y_w, y_l) \sim \mathcal{D}} \bigg[\underbrace{\sigma(\hat{r}_\theta(x, y_l) - \hat{r}_\theta (x, y_w))}_\text{larger values if reward order is incorrect}\bigg[\underbrace{\nabla_\theta\log \pi(y_w \mid x)}_\text{encourage $y_w$} - \underbrace{\nabla_\theta\log\pi(y_l \mid x)}_\text{discourage $y_l$}\bigg]\bigg],
\end{multline*}

where the implicit reward is expressed as $\hat{r}_\theta(x, y) = \beta \log \frac{\pi_\theta(y \mid x)}{\pi_\text{ref}(y \mid x)}$, utilizing both the learned language model $\pi_\theta$ and a reference policy $\pi_\text{ref}$ (additional details are presented in Section~\ref{sec:theory}). Conceptually, this gradient acts to promote the probability of preferred completions $y_w$ while suppressing the probability of less-preferred completions $y_l$. Critically, each sample is weighted by the extent to which the implicit reward function $\hat{r}_\theta$ mistakenly rates $y_l$ higher than $y_w$, scaled by $\beta$, thereby modulating the influence in accordance with the KL constraint's strength. Our empirical findings highlight the necessity of this weighting: omitting it, as in a simple variant of the method, leads to poor language model performance, including degeneration issues (see Appendix Table~\ref{tab:unlikelihood_generations}).

\textbf{DPO overview.} The standard procedure for DPO consists of: 1) Generating completions $y_1, y_2 \sim \pi_\text{ref}(\cdot \mid x)$ for each input $x$, assigning human preference labels, and forming an offline preference dataset $\mathcal{D} = \{x^{(i)}, y_w^{(i)}, y_l^{(i)}\}_{i=1}^N$; and 2) updating the language model $\pi_\theta$ to minimize $\mathcal{L}_\text{DPO}$, conditioned on $\pi_\text{ref}$, $\mathcal{D}$, and a chosen $\beta$. 
In practical applications, practitioners often leverage existing, public preference datasets instead of generating fresh samples and collecting new human preference data. Given that these datasets are typically drawn using $\pi^\text{SFT}$, we adopt $\pi_\text{ref} = \pi^\text{SFT}$ as the initial reference policy when it is accessible. If $\pi^\text{SFT}$ is unavailable, $\pi_\text{ref}$ is instead initialized by maximizing the likelihood over preferred completions $(x, y_w)$, specifically: ${\pi_\text{ref} = \argmax_{\pi}\mathbb{E}_{x, y_w \sim \mathcal{D}}\left[\log \pi(y_w \mid x)\right]}$. This approach is intended to reduce distributional discrepancies between the inaccessible true reference distribution and the $\pi_\text{ref}$ used within DPO. For more implementation specifics and hyperparameter choices, refer to Appendix~\ref{app:implementation}.

\section{Theoretical Analysis of DPO} This section develops a deeper theoretical understanding of the DPO algorithm, formalizes supporting results, and discusses the strengths of DPO as contrasted with common actor-critic algorithms in RLHF, such as PPO~\cite{schulman2017proximal}.

\label{sec:theory}

\subsection{Your Language Model Is Secretly a Reward Model} DPO replaces both reward modeling and RL with a unified maximum likelihood approach. The loss function in Eq.~\ref{eq:main_eq} coincides with the Bradley-Terry model using a reward defined by $r^*(x, y) = \beta \log\frac{\pi^*_\theta(y \mid x)}{\pi_\text{ref}(y \mid x)}$. Here, we optimize the parameterized model $\pi_{\theta}$ in a fashion analogous to direct reward model training in Eq.~\ref{eq:reward_model}, given a suitable change of variables. This section constructs the theoretical justification for this reparameterization, demonstrating that it preserves the generality of representable reward models and enables exact retrieval of the optimal policy. We first introduce a formal equivalence notion for reward functions.

\begin{definition}
Two reward functions $r(x, y)$ and $r'(x, y)$ are said to be equivalent if and only if there exists a function $f$ such that ${r(x, y)-r'(x, y) = f(x)}$.
\end{definition}

This definition establishes an equivalence relation that segments the set of reward functions into distinct equivalence classes. From this, we can derive two immediate lemmas:

\begin{lemma}\label{lemma:same_prefrence}
Within the context of the Plackett-Luce preference framework—encompassing the Bradley-Terry model—any pair of reward functions belonging to the same equivalence class result in identical preference distributions.
\end{lemma}

\begin{lemma}\label{lemma:same_policy}
Any two reward functions from the same equivalence class yield the same optimal policy for the constrained reinforcement learning problem.
\end{lemma}

The proofs for these lemmas are direct and can be found in Appendix \ref{app:lemma1}. The first lemma echoes a widely recognized issue of under-specification present in models from the Plackett-Luce family \cite{plackett1975analysis}. This ambiguity often necessitates additional identifiability constraints to ensure meaningful maximum likelihood estimates as formulated in Eq. \ref{eq:reward_model} \cite{bong2022generalized}. The second lemma formalizes that within a given equivalence class, all reward functions are optimal with respect to the same policy; thus, for the purposes of our main objective, it suffices to recover any representative from the optimal equivalence class. We formally state the following theorem, whose proof appears in Appendix~\ref{app:thm1}:

\begin{theorem}\label{thm:main}
Given mild conditions, any reward equivalence class consistent with Plackett-Luce models (including Bradley-Terry as a specific case) can be parameterized by ${r(x, y) = \beta \log \frac{\pi(y\mid x)}{\pi_\text{ref}(y\mid x)}}$, where $\pi(y\mid x)$ is an arbitrary model and $\pi_\text{ref}(y \mid x)$ serves as a fixed reference model.
\end{theorem}

\begin{sproof}
Suppose we have an arbitrary reward function $r(x, y)$ that determines an associated optimal model $\pi_r(y \mid x)$, as defined in Eq. \ref{eq:op_policy}. We demonstrate that one can represent a reward function in the equivalence class of $r$ by employing the suggested reparameterization. Consider the projection operator $f$ defined as
\begin{equation}
    f(r; \pi_\text{ref}, \beta)(x, y) = r(x, y) - \beta\log\sum_{y}\pi_\text{ref}(y\mid x)\exp\left(\frac{1}{\beta}r(x, y)\right)
\end{equation}
This operator $f$ serves to normalize the reward function by subtracting the log partition function (depending only on $x$) for $\pi_r$. Since this normalization term depends solely on $x$, $f(r; \pi_\text{ref}, \beta)(x, y)$ remains within the equivalence class of $r(x, y)$. By substituting the right-hand side of Eq.~\ref{eq:main_eq} (applicable for any $r$), we have $f(r; \pi_\text{ref}, \beta)(x, y) = \beta \log \frac{\pi_r(y\mid x)}{\pi_\text{ref}(y\mid x)}$. Thus, the mapping $f$ returns an element from the equivalence class of $r$ in the desired log-ratio form, and the proposed reparameterization entails no loss of generality in the model class.
\end{sproof}

An alternative interpretation of Theorem~\ref{thm:main} is that it precisely determines which specific reward function is chosen by the DPO reparameterization within each equivalence class—namely, the one that fulfills the following condition:

\begin{equation}\label{eq:lag_p}
     \sum_{y}\underbrace{\pi_\text{ref}(y\mid x)\exp\left(\frac{1}{\beta}r(x, y)\right)}_{=\pi(y\mid x)\text{, using Thm.~\ref{thm:main} reparam.}} = 1,
\end{equation}

That is, $\pi(y\mid x)$ indeed defines a legitimate probability distribution whose elements are non-negative and collectively sum to one. Moreover, as highlighted by Eq.~\ref{eq:op_policy}, Eq.~\ref{eq:lag_p} corresponds to the partition function for the optimal policy defined by the reward $r(x, y)$. The central innovation in the DPO algorithm is that, by introducing certain constraints on the otherwise unconstrained Plackett-Luce (and specifically the Bradley-Terry) class of preference models, one can both retain the richness of the class of permissible reward functions and ensure that the optimal policy from Eq.~\ref{eq:op_policy} remains analytically computable for any prompt $x$.

\subsection{Instability of Actor-Critic Algorithms}
Our analytical framework is also useful for uncovering sources of instability in commonly used actor-critic algorithms for RLHF, such as PPO. We adopt the RLHF methodology, with particular attention to the RL fine-tuning phase described in Section \ref{section:prelims}. By drawing parallels with the control as inference paradigm \cite{levine2018reinforcement}, applied to the constrained RL setting as formalized in \ref{eq:RL}, we consider a parameterized policy $\pi_{\theta}(y\mid x)$. The training procedure involves minimizing $\mathbb{D}_{\text{KL}}[\pi_{\theta}(y|x) \mid \mid \pi^*(y\mid x)]$, where $\pi^*$ denotes the optimal policy obtained from Eq. \ref{eq:optimum_model} as a function of the reward $r_{\phi}(y, x)$. After some manipulation, this objective leads to the following optimization problem:

\begin{equation}\label{eq:AC}
    \max_{\pi_{\theta}}\mathbb{E}_{\pi_{\theta}(y\mid x)}\bigg[\underbrace{r_{\phi}(x, y) -\beta\log\sum_{y}\pi_\text{ref}(y\mid x)\exp\left(\frac{1}{\beta}r_{\phi}(x, y)\right)}_{f(r_{\phi}, \pi_\text{ref}, \beta)} - \underbrace{\beta\log\frac{\pi_{\theta}(y\mid x)}{\pi_\text{ref}(y\mid x)}}_{\text{KL}}\bigg]
\end{equation}

The same objective function has been optimized in earlier studies 
\citep{ziegler2020finetuning, stiennon2022learning, bai2022training, ouyang2022training} by employing a DPO-equivalent reward formulation for the class $r_{\phi}$. Within this context, the normalization component in $f(r_{\phi}, \pi_\text{ref}, \beta)$ can be viewed as the soft value function associated with the reference policy $\pi_\text{ref}$. Although this normalization term does not alter the optimal policy, omitting it can result in high-variance policy gradients, which may hinder the stability of the learning process. One option is to address this normalization by leveraging a learned value function, although this approach can present optimization challenges as well. Another approach, found in previous literature, involves normalizing rewards against a human-generated baseline, effectively a Monte-Carlo estimate of the normalization term based on a single sample. By contrast, the DPO reparameterization introduces a reward function that circumvents the need for such baselines.

\section{Experiments}
This section presents an empirical assessment of DPO's effectiveness in training policies using preference data. Initially, within a controlled text generation framework, we examine the efficiency of DPO in balancing reward maximization with minimization of KL-divergence from the reference policy, comparing its performance against prevalent preference learning algorithms such as PPO. Subsequently, DPO's capabilities are tested on larger-scale models and more complex RLHF scenarios, including tasks in summarization and conversational dialogue. Our results indicate that, even with minimal hyperparameter tuning, DPO typically matches or surpasses established baselines such as PPO-based RLHF and the strategy of selecting the best trajectory out of $N$ samples according to a learned reward. Prior to detailing these outcomes, we outline the experimental methodology; further specifics can be found in Appendix~\ref{app:exp_details}.

\textbf{Tasks.} We conduct experiments across three distinct open-ended text generation domains. Across all settings, the algorithms aim to acquire a policy using a preference dataset $\mathcal{D}=\bigl\{x^{(i)}, y_w^{(i)}, y_l^{(i)}\bigr\}_{i=1}^N$. In the case of \textbf{controlled sentiment generation}, $x$ represents the initial segment of a movie review sourced from the IMDb dataset \cite{maas-EtAl:2011:ACL-HLT2011}, with the policy required to produce $y$ exhibiting positive sentiment. To facilitate systematic evaluation, we synthesize preference pairs between generations via a pre-trained sentiment classifier, enforcing $p(\text{positive}\mid x,y_w)>p(\text{positive}\mid x,y_l)$. For SFT, GPT-2-large is fine-tuned on training-set reviews from the IMDB corpus until performance plateaus (see App~\ref{app:sentiment_details} for more details). In the \textbf{summarization} task, $x$ corresponds to a Reddit forum post, and the policy's objective is to generate a summary $y$ highlighting the central ideas. Following established protocols, we employ the Reddit TL;DR summarization dataset \citep{volske-etal-2017-tl} alongside a human preference dataset collected by \citeauthor{stiennon2022learning}. An SFT model is used, fine-tuned on human-authored summaries of forum posts\footnote{\url{https://huggingface.co/CarperAI/openai_summarize_tldr_sft}}, utilizing the TRLX \citep{leandro_von_werra_2023_7790115} RLHF toolkit. The dataset of human preferences is built from model samples distinct from—but similarly trained to—the primary SFT model, as reported by \citeauthor{stiennon2022learning}. The final domain is \textbf{single-turn dialogue}, where $x$ denotes a user's query, varying from technical inquiries (e.g., astrophysics) to personal advice. Here, the policy is expected to generate $y$, a response that is both informative and engaging. We utilize the Anthropic Helpful and Harmless dialogue dataset \citep{bai2022training}, which comprises 170k dialogues featuring a human and an automated assistant. Each dialogue concludes with two model-generated responses, accompanied by an annotation indicating which response a human rater favored. As a suitable pre-trained SFT model is lacking in this scenario, we construct an SFT model by fine-tuning a generic language model solely on preferred responses.

\textbf{Evaluation.} We assess our models using two distinct evaluation strategies. For a precise examination of each algorithm’s capability to optimize the constrained reward maximization objective, we conduct controlled sentiment generation experiments, wherein algorithms are evaluated based on their reward versus KL-divergence trade-off relative to the reference policy. This analysis is feasible due to our access to the exact reward function (a sentiment classifier). However, in realistic applications where the ground-truth reward is unavailable, we compare algorithms by measuring their \textit{win rate} against a baseline, leveraging GPT-4 as a surrogate for human judgment to evaluate summary quality and dialogue helpfulness in summarization and single-turn dialogue tasks. For summarization, the reference summaries in the test data serve as the baseline, while in dialogue, we use the test set’s preferred response. Despite prior findings that language models outperform traditional metrics as automated evaluators \citep{Chen2023ExploringTU}, we additionally conduct a human study, described in Sec.~\ref{sec:human-judgments}, to support the validity of using GPT-4 for our evaluations. Our results show a strong alignment between GPT-4 and human raters, with human-GPT-4 agreement typically matching or surpassing agreement rates between human annotators.

\textbf{Methods.} Besides DPO, we benchmark multiple established techniques for aligning language models with human preferences. For summarization, we utilize zero-shot prompting with \textbf{GPT-J} \citep{gpt-j}, and for dialogue, we apply 2-shot prompting with \textbf{Pythia-2.8B} \citep{biderman2023pythia}. Our evaluations also include the \textbf{SFT} model and the \textbf{Preferred-FT} model, the latter being further trained through supervised fine-tuning on preferred completions $y_w$ generated either by SFT (in controlled sentiment and summarization tasks) or a generic LM (in dialogue tasks). We additionally assess the \textbf{Unlikelihood} method~\citep{welleck2019neural}, which modifies the policy to increase the probability assigned to $y_w$ while suppressing that of $y_l$, controlled by a tunable parameter $\alpha\in[0,1]$ applied to the unlikelihood term. Our comparisons further include \textbf{PPO} \citep{schulman2017proximal} using a reward model trained from preference data, as well as \textbf{PPO-GT}, an oracle method trained with access to the ground-truth reward available in the controlled sentiment scenario. For sentiment experiments, we deploy two PPO-GT variants: a standard off-the-shelf version \cite{leandro_von_werra_2023_7790115} and a customized version that normalizes rewards and tunes hyperparameters for better outcomes (these enhancements are likewise employed for standard PPO with learned reward models). Finally, we test the \textbf{Best of $N$} strategy, which generates $N$ responses from the SFT (or Preferred-FT in dialogue), selecting the highest-rewarded output as per a reward model trained on preference data. Although this method isolates the impact of the reward model from PPO’s optimization process, it is computationally prohibitive at test time even for moderate $N$, given the necessity to sample $N$ completions per input.

\subsection{How well can DPO optimize the RLHF objective?}

Typical RLHF methods employ a KL-regularized reward maximization objective, striking a balance between maximizing reward and keeping the learned policy close to the reference policy. As a result, proper evaluation of different algorithms requires jointly considering both the attained reward and the KL divergence; increasing reward at the expense of a substantial KL increase is often undesirable. Figure~\ref{fig:frontier-tldr-main} illustrates the reward versus KL trade-off curves for various approaches on a sentiment analysis task. For each algorithm, we perform multiple training sessions, varying a single hyperparameter controlling policy conservativeness across runs (for PPO, we sweep the target KL over $\{3,6,9,12\}$; for DPO, $\beta$ over $\{0.05,0.1,1,5\}$; for unlikelihood, $\alpha$ over $\{0.05,0.1,0.5,1\}$; for preferred-FT, we vary the random seed). This comprehensive search encompasses a total of 22 experiments. Following every 100 training steps up to convergence, policies are assessed using a set of test prompts, with both the mean reward (using the actual reward function) and the mean sequence-level KL\footnote{Namely, the sum across all timesteps of the KL divergences.} with respect to the reference policy $\text{KL}\left(\pi\mid \mid \pi_\text{ref}\right)$ reported. Our results show that DPO consistently constructs a more favorable frontier than the other approaches, yielding superior rewards while keeping KL low. These findings are especially remarkable for several reasons. Primarily, although DPO and PPO are designed to optimize the same underlying objective, DPO demonstrates considerably greater efficiency, with its reward/KL curve strictly outperforming PPO. Furthermore, DPO’s frontier remains better than that of PPO even in the scenario where PPO is given access to ground-truth rewards (PPO-GT).

\subsection{Can DPO scale to real preference datasets?}
\label{sec:dpo-real-datasets}
We proceed to assess DPO’s ability to fine-tune models on tasks such as summarization and single-turn dialogue. In the context of summarization, existing automatic metrics like ROUGE have been shown to correlate poorly with human preference~\citep{stiennon2022learning}. Earlier studies have demonstrated that leveraging PPO to fine-tune language models on human preference data leads to more useful summaries. To compare methods, we generate outputs on the test portion of the TL;DR summarization dataset and calculate the average win rate of these outputs when compared to reference completions from the test set. Samples from each method are generated across a range of temperatures from 0.0 to 1.0; corresponding win rates are visualized in Figure~\ref{fig:frontier-tldr-main} (right). For these experiments, DPO, PPO, and Preferred-FT are all initialized from the same GPT-J SFT checkpoint\footnote{\url{https://huggingface.co/CarperAI/openai_summarize_tldr_sft}}. Our results indicate that DPO achieves an approximate win rate of 61\% at a temperature of 0.0, outperforming PPO, which reaches about 57\% at its optimal temperature (also 0.0). DPO attains a superior maximal win rate in comparison to the best-of-$N$ baseline. We acknowledge that DPO’s $\beta$ parameter was not systematically optimized, and as such, the results might not fully reflect the method’s capabilities. Additionally, DPO exhibits substantially greater resilience to variations in sampling temperature, whereas PPO’s performance can decline to that of the initial GPT-J model when higher temperatures are used. In contrast, Preferred-FT does not offer noticeable gains over the SFT baseline. A direct comparison between DPO and PPO via human evaluation is presented in Section~\ref{sec:human-judgments}; there, samples from DPO at a temperature of 0.25 were preferred by annotators 58\% of the time compared to PPO samples at the same temperature.

For single-turn dialogue evaluation, we analyze several approaches using the subset of the Anthropic HH dataset’s test split \citep{bai2022training} that consists of single-round interactions between a human and an assistant. In our GPT-4-based assessments, we reference the preferred test completions to determine win rates across methods. Since there is no established SFT baseline for this setting, our workflow initiates with the pre-trained Pythia-2.8B model. We employ Preferred-FT to construct a reference model, training it on selected completions to maintain distributional similarity, before further training via DPO. Our experiments also include comparisons with the best outcome among 128 Preferred-FT completions—having observed that the Best of $N$ approach levels off at 128 completions for this benchmark (see Appendix Figure~\ref{fig:best-of-n})—as well as a two-shot prompted variant of the Pythia-2.8B base model. The results show DPO achieves performance equal to or exceeding that of the other methods when the optimal temperature is chosen. Additionally, we assess an RLHF model trained with PPO on the Anthropic HH dataset \footnote{\url{https://huggingface.co/reciprocate/ppo_hh_pythia-6B}}, utilizing a reputable implementation \footnote{\url{https://github.com/CarperAI/trlx/tree/main/examples/hh}}; however, no combination of prompt or sampling temperature led this model to surpass the baseline Pythia-2.8B model. Considering insights from the TL;DR task and recognizing that both PPO and Best of 128 seek to optimize the same reward, we regard Best of 128 as an approximate stand-in for PPO-level performance. In summary, DPO uniquely stands out as a computationally efficient strategy that consistently outperforms the preferred completions on the Anthropic HH dataset and delivers results comparable to or better than the more resource-intensive Best of 128 baseline. Figure~\ref{fig:dialogue-main} further illustrates that DPO attains peak performance after relatively few training steps.

\subsection{Generalization to a new input distribution}

To more rigorously assess how PPO and DPO handle distributional changes, we examine their generalization performance by evaluating the policies, initially trained on Reddit TL;DR summarization, on a distinct dataset: the test split of the CNN/DailyMail corpus \citep{nallapati-etal-2016-abstractive}. We utilize the optimal sampling temperatures identified from the TL;DR task (0 and 0.25) during this evaluation. The corresponding outcomes are displayed in Table~\ref{tab:ood}. To determine the win rates, we used GPT-4 to compare generated summaries against the dataset’s gold-standard references, employing the same GPT-4 (C) prompt as in the Reddit TL;DR setting, but substituting “forum post” with “news article”. Across this distribution shift, DPO remains superior to PPO by a notable margin. These findings indicate that DPO can generalize at least as effectively as PPO, even without the benefit of the extra unlabeled Reddit TL;DR prompts utilized by PPO.

\subsection{Validating GPT-4 judgments with human judgments}
\label{sec:human-judgments}
To assess the trustworthiness of GPT-4-based evaluations, we conduct a user study using outputs from the TL;DR summarization experiments and two different GPT-4 evaluation prompts. The \textbf{GPT-4 (S)} (simple) prompt requests a judgment on which summary best captures the key content of a post. The \textbf{GPT-4 (C)} (concise) prompt further asks which summary is more succinct; we include this prompt after observing that, under the \textbf{GPT-4 (S)} prompt, GPT-4 tends to prefer lengthier and more redundant summaries than human annotators do. The full prompt wording can be found in Appendix~\ref{app:prompts}. We execute three sets of comparisons: the top-performing model (DPO, temp. 0.25), the lowest-performing (PPO, temp. 1.0), and an intermediate method (SFT, temp. 0.25), each assessed against PPO with greedy sampling at its best temperature. This approach samples across varying quality levels. Our findings indicate that, for both prompts, GPT-4’s preferences are as consistent with those of human evaluators as the agreement rate among the humans themselves (note that, due to the number of available annotators, multiple human ratings are collected only for the DPO and PPO-1 comparisons). Importantly, the \textbf{GPT-4 (C)} prompt produces win rates that more closely reflect human judgments, and as such, we rely on this prompt for our main evaluation in Section~\ref{sec:dpo-real-datasets}. Details about the human study design, interface, and a list of participants are available in Appendix~\ref{app:human-study}.

\section{Discussion}
Preference-based learning offers an effective and scalable approach for aligning and enhancing the capabilities of language models. In this work, we have presented DPO—a straightforward method for training language models on preference data, circumventing the need for reinforcement learning altogether. DPO departs from the common practice of reframing preference learning within a reinforcement learning context; instead, it establishes a direct relationship between policy outputs and reward functions. This permits direct optimization to human preferences using a cross-entropy loss function, entirely avoiding reinforcement learning algorithms while maintaining full generality. DPO matches or outperforms prominent RLHF strategies, such as those relying on PPO, often without extensive hyperparameter tuning. Consequently, DPO streamlines the process of leveraging human preferences for training language models.

\textbf{Limitations \& Future Work.} Our findings highlight several avenues that warrant further investigation. One open question concerns the out-of-distribution generalization capabilities of the DPO-trained policy compared to those models optimized with explicit reward signals; preliminary comparisons indicate comparable generalization to PPO-based approaches, but more systematic evaluations are necessary. Another consideration is whether DPO’s self-labeling methods can effectively exploit unannotated prompt data. We are also interested in characterizing how reward over-optimization might arise specifically in direct preference optimization, and whether the observed minor decline in Figure~\ref{fig:dialogue-main}-right is symptomatic of this issue. Although our experiments focus on models up to 6B parameters, scaling DPO to the size of leading-edge architectures presents a promising research direction. In terms of evaluation methodology, we note that GPT-4’s estimated win rates exhibit sensitivity to prompt formulations, pointing to a need for better strategies for extracting reliable automated assessments. Ultimately, DPO’s utility extends beyond language modeling from human preferences, with potential for broader application in the training of generative models across diverse modalities.